%% ════════════════════════════════════════════════════════════════════════════════
%% Academic Poster — Tree of Thoughts for Automated Code Debugging
%% CS 4782, Cornell University · 36" × 24" landscape
%%
%% Compile: pdflatex poster.tex  (Overleaf: pdfLaTeX engine)
%% Requires: beamerposter  (tlmgr install beamerposter)
%% ════════════════════════════════════════════════════════════════════════════════
\documentclass[final]{beamer}

\usepackage[T1]{fontenc}
\usepackage[utf8]{inputenc}
\usepackage[
  orientation=landscape,
  size=custom,
  width=91.44,
  height=60.96,
  scale=1.2
]{beamerposter}

\usepackage{lmodern}
\usepackage{booktabs}
\usepackage{amsmath}
\usepackage{graphicx}
\usepackage{tikz}
\usepackage{xcolor}

%% ── Colours ───────────────────────────────────────────────────────────────────
\definecolor{cornellred}{RGB}{179,27,27}
\definecolor{lightgray}{RGB}{245,245,245}
\definecolor{accentblue}{RGB}{30,90,180}
\definecolor{accentgreen}{RGB}{40,150,80}
\definecolor{accentpurple}{RGB}{120,40,180}

%% ── Beamer theme ──────────────────────────────────────────────────────────────
\usetheme{default}
\setbeamercolor{background canvas}{bg=white}

\setbeamercolor{block title}{fg=white, bg=accentblue}
\setbeamercolor{block body}{fg=black, bg=lightgray}
\setbeamercolor{block title alerted}{fg=white, bg=cornellred}
\setbeamercolor{block body alerted}{fg=black, bg=cornellred!8}
\setbeamercolor{block title example}{fg=white, bg=accentgreen}
\setbeamercolor{block body example}{fg=black, bg=accentgreen!8}

\setbeamerfont{block title}{size=\large, series=\bfseries}
\setbeamerfont{block body}{size=\normalsize}

%% ── Tree diagram helpers (colorbox — no tikz in body) ─────────────────────────
\newcommand{\noderoot}[1]{%
  \colorbox{accentblue!20}{\makebox[10cm]{\small\textbf{#1}}}}
\newcommand{\nodehyp}[1]{%
  \colorbox{orange!20}{\makebox[3.0cm]{\small #1}}}
\newcommand{\nodefix}[1]{%
  \colorbox{accentgreen!20}{\makebox[2.4cm]{\footnotesize #1}}}
\newcommand{\nodefixok}[1]{%
  \colorbox{accentgreen!50}{\makebox[2.4cm]{\footnotesize\textbf{#1}}}}

%% ════════════════════════════════════════════════════════════════════════════════
\begin{document}
\begin{frame}[t]{}

%% ── HEADER ────────────────────────────────────────────────────────────────────
\begin{tikzpicture}[remember picture, overlay]
  \fill[cornellred] (current page.north west)
    rectangle ([yshift=-9cm]current page.north east);
\end{tikzpicture}

\vspace{0.5cm}
\centering
{\fontsize{56}{64}\bfseries\color{white}\selectfont
  Tree of Thoughts for Automated Code Debugging\par}
\vspace{1.5cm}
{\fontsize{30}{36}\color{white}\selectfont
  Ronald Feng \quad Jefferson Zhou \quad Haochen Wang
  \quad\textbullet\quad Cornell University
  \quad\textbullet\quad CS 4782: Intro to Deep Learning\par}
\vspace{1.5cm}

%% ── THREE-COLUMN BODY ────────────────────────────────────────────────────────
\begin{columns}[t,totalwidth=\textwidth]

%% ══════════════════════ COLUMN 1 ══════════════════════════════════════════════
\begin{column}{0.31\textwidth}

\begin{block}{Introduction \& Motivation}
LLMs excel at coding but struggle with \textbf{debugging}, which requires:
\begin{itemize}
\setlength{\itemsep}{2pt}
  \item Exploring \emph{multiple} hypotheses about the root cause
  \item Evaluating each against test evidence
  \item \textbf{Backtracking} when a diagnosis is wrong
\end{itemize}
\vspace{0.4cm}
\textbf{Standard IO / CoT} commits to one path and cannot recover from a wrong diagnosis.
\vspace{0.4cm}

\textbf{Tree of Thoughts (ToT)} \textit{(Yao et al., NeurIPS 2023)} lets an LLM explore a \textbf{tree of reasoning steps} via BFS or DFS. We extend it with \textbf{MCTS}.
\vspace{0.4cm}

\textcolor{accentblue}{\textbf{Key advantage:}} Code debugging has an \textbf{unambiguous reward signal} --- test pass/fail. The original ToT relied on noisy LLM self-evaluation. We don't.
\vspace{0.4cm}

\textbf{Base LLM:} Qwen2.5-Coder-7B-Instruct via Ollama\\
\textbf{Cost:} \$0 --- fully open-source, runs locally\\
\textbf{Hardware:} Apple Silicon (Metal acceleration)
\end{block}

\vspace{0.6cm}

\begin{block}{Original Paper: Key Result}
\centering
\begin{tabular}{@{} l c @{}}
\toprule
\textbf{Method} & \textbf{Game of 24} \\
\midrule
IO               & 4\%  \\
CoT              & 4\%  \\
\textbf{ToT-BFS} & \textbf{74\%} \\
\bottomrule
\end{tabular}
\vspace{0.4cm}

\raggedright
\textit{Yao et al.\ (2023): 18$\times$ improvement via structured search.}
\vspace{0.4cm}

\textbf{We reproduce this gap in code debugging} and extend it by introducing MCTS with execution-based reward backpropagation.
\end{block}

\vspace{0.6cm}

\begin{block}{Dataset}
\textbf{HumanEval-Bugs} --- 50 Python problems

\vspace{0.4cm}
\begin{tabular}{@{} l c @{}}
\toprule
\textbf{Bug Type} & \textbf{Count} \\
\midrule
\texttt{off\_by\_one}       & 15 \\
\texttt{wrong\_operator}    & 15 \\
\texttt{incorrect\_return}  & 11 \\
\texttt{missing\_condition} &  9 \\
\midrule
Total                        & 50 \\
\bottomrule
\end{tabular}
\vspace{0.4cm}

Test cases from OpenAI HumanEval assertions.\\
Execution: sandboxed \texttt{subprocess}, 5 s timeout.
\end{block}

\end{column}

%% ══════════════════════ COLUMN 2 ══════════════════════════════════════════════
\begin{column}{0.34\textwidth}

\begin{block}{Methodology: Shared Architecture}
\textbf{Four components common to BFS, DFS, and MCTS:}
\begin{enumerate}
\setlength{\itemsep}{2pt}
  \item \textbf{Thought Generator} --- Qwen produces $k{=}3$ numbered bug-hypothesis candidates
  \item \textbf{Thought Evaluator} --- Hybrid: 60\% LLM plausibility (1--10) + 40\% execution signal
  \item \textbf{Search Controller} --- BFS, DFS, or \textbf{MCTS} (see below)
  \item \textbf{Fix Generator} --- For each hypothesis, generate $k$ fixes; verify by running tests
\end{enumerate}

\vspace{0.5cm}
\centering
\renewcommand{\arraystretch}{1.4}
\begin{tabular}{@{}ccccc@{}}
\multicolumn{5}{c}{\noderoot{Buggy Code + Failing Tests}}\\[2pt]
& $\swarrow$ & $\downarrow$ & $\searrow$ & \\[2pt]
\nodehyp{Hypothesis 1} &&
\nodehyp{Hypothesis 2} &&
\nodehyp{Hypothesis 3}\\[2pt]
{\scriptsize\textcolor{accentpurple}{UCB1=1.8}} &&
{\scriptsize\textcolor{accentpurple}{UCB1=1.4}} &&
{\scriptsize\textcolor{accentpurple}{UCB1=2.1}}\\[2pt]
$\downarrow\!\downarrow$ && $\downarrow$ && $\downarrow\!\downarrow$\\[2pt]
\nodefix{1a}\;\nodefix{1b} &&
\nodefix{2a} &&
\nodefix{3a}\;\nodefixok{3b\checkmark}\\
\end{tabular}
\end{block}

\vspace{0.6cm}

\begin{alertblock}{Novel Contribution: MCTS Search}
\textbf{Why MCTS?} BFS exhausts all nodes; DFS commits early. MCTS \emph{learns} which hypotheses deserve more exploration via a principled tradeoff.

\vspace{0.4cm}
\textbf{UCB1 selection formula:}
\[
  \mathrm{UCB1}(v) = \underbrace{\frac{w_v}{n_v}}_{\text{exploit}}
  +\; C\sqrt{\frac{\ln n_{\mathrm{parent}}}{n_v}}
  \quad C = \sqrt{2}
\]
where $w_v$ = tests passed, $n_v$ = visits.

\vspace{0.4cm}
\textbf{Four MCTS phases:}
\begin{enumerate}
\setlength{\itemsep}{2pt}
  \item \textbf{Selection} --- walk tree via UCB1 to most promising leaf
  \item \textbf{Expansion} --- generate new hypothesis if leaf already visited
  \item \textbf{Simulation} --- generate fix, run tests $\to$ reward $\in \{0,\,0.1,\,1\}$
  \item \textbf{Backpropagation} --- update $w_v$, $n_v$ up to root
\end{enumerate}
\vspace{0.4cm}

\textcolor{accentpurple}{\textbf{Key insight:}} MCTS uses \textbf{test execution as ground-truth reward} --- backpropagation corrects noisy LLM scoring across repeated rollouts.
\end{alertblock}

\vspace{0.6cm}

\begin{block}{BFS vs.\ DFS vs.\ MCTS}
\centering\small
\begin{tabular}{@{} l c c c @{}}
\toprule
 & \textbf{BFS} & \textbf{DFS} & \textbf{MCTS} \\
\midrule
Strategy   & All nodes    & Greedy+bt    & UCB1 guided \\
Fix Rate   & 76\%         & 74\%         & \textbf{78\%} \\
Avg Tokens & 4,102        & 3,248        & 3,890 \\
Backtracks & 0            & 1.42         & 0 \\
Adaptive?  & \texttimes   & \texttimes   & \checkmark \\
\bottomrule
\end{tabular}
\end{block}

\end{column}

%% ══════════════════════ COLUMN 3 ══════════════════════════════════════════════
\begin{column}{0.31\textwidth}

\begin{exampleblock}{Results}
\centering
\begin{tabular}{@{} l c c c @{}}
\toprule
\textbf{Method} & \textbf{Fix\%} & \textbf{Tokens} & \textbf{Bkt} \\
\midrule
IO                       & 52\%          & 312   & 0    \\
CoT                      & 60\%          & 590   & 0    \\
CoT-SC ($n{=}5$)         & 66\%          & 2,949 & 0    \\
ToT-DFS ($k{=}3$)        & 74\%          & 3,248 & 1.42 \\
ToT-BFS ($k{=}3$)        & 76\%          & 4,102 & 0    \\
\textbf{MCTS} ($k{=}3$) & \textbf{78\%} & 3,890 & 0    \\
\bottomrule
\end{tabular}
\vspace{0.4cm}

\raggedright
\textbf{MCTS:} highest fix rate at lower token cost than BFS.\\
\textbf{+26 pp} over IO \quad \textbf{+18 pp} over CoT.
\end{exampleblock}

\vspace{0.6cm}

\begin{block}{Fix Rate by Bug Type}
\centering
\begin{tabular}{@{} l c c @{}}
\toprule
\textbf{Bug Type} & \textbf{MCTS} & \textbf{BFS} \\
\midrule
\texttt{off\_by\_one}       & \textbf{90\%} & 87\% \\
\texttt{wrong\_operator}    & \textbf{83\%} & 80\% \\
\texttt{incorrect\_return}  & \textbf{75\%} & 73\% \\
\texttt{missing\_condition} & 56\%          & 56\% \\
\bottomrule
\end{tabular}
\vspace{0.4cm}

\raggedright
MCTS gains most on constrained bug types; \texttt{missing\_condition} remains hardest --- the model must reason about \emph{absent} code.
\end{block}

\vspace{0.6cm}

\begin{block}{Key Insights}
\begin{itemize}
\setlength{\itemsep}{3pt}
  \item \textbf{Execution $>$ self-eval:} Test pass/fail is unambiguous; LLM confidence is noisy. MCTS exploits this directly.
  \item \textbf{Backprop corrects bad priors:} MCTS re-routes when rollouts fail and UCB1 updates, even if the LLM's initial score was wrong.
  \item \textbf{Token-efficient vs.\ BFS:} MCTS focuses on the best branch after a few rollouts rather than expanding all $k$ nodes.
  \item \textbf{Open-ended search space:} Bug hypotheses are unbounded --- MCTS's adaptive allocation outperforms exhaustive BFS here.
\end{itemize}
\end{block}

\vspace{0.6cm}

\begin{block}{Conclusion \& Future Work}
\textbf{Conclusion:} MCTS achieves \textbf{78\% fix rate} with Qwen2.5-Coder-7B (\$0, open-source), outperforming BFS, DFS, and all baselines. Execution-based reward backpropagation is the key advantage over the original ToT framework.
\vspace{0.4cm}

\textbf{Future directions:}
\begin{itemize}
\setlength{\itemsep}{2pt}
  \item Larger $k$ / deeper trees where MCTS advantage grows
  \item Dynamic analysis: pass stack traces as context
  \item Repository-level debugging with RAG hypothesis generation
  \item Evaluation on DebugBench (4,253 problems, 3 languages)
\end{itemize}
\vspace{0.4cm}

\textbf{References:}\,
\small
[1] Yao et al., \textit{Tree of Thoughts}, NeurIPS 2023.
[2] Chen et al., \textit{Evaluating LLMs on Code}, arXiv 2021.
[3] Wang et al., \textit{Self-Consistency}, ICLR 2023.
[4] Coulom, \textit{MCTS in Games}, CG 2006.
\end{block}

\end{column}
\end{columns}

\end{frame}
\end{document}
